\problemname{Björne's Store}
As is well known, bears hibernate during the winter. But at the equator it does not get dark and cold during the winter,
so ever since a bear family moved south many years ago, the bears in the bear community that arose
have started hibernating at completely different times of the year.
Björne has opened a store and is worried about thieves during the part of the year when he is asleep.
He therefore needs to hire other bears so that there is always someone employed who is awake and can guard the store.

In total, $N$ bears have applied for the job.
Each bear goes into hibernation on a certain day of the year to sleep for exactly $d$ consecutive days, and cannot work during any of these days.
How many bears does Björne need to hire at minimum, given that Björne will guard his store himself when he is awake?

Assume that every year has $365$ days (bears living at the equator do not care about leap years), and that:
\begin{enumerate}
  \item January has 31 days,
  \item February has 28 days,
  \item March has 31 days,
  \item April has 30 days,
  \item May has 31 days,
  \item June has 30 days,
  \item July has 31 days,
  \item August has 31 days,
  \item September has 30 days,
  \item October has 31 days,
  \item November has 30 days, and
  \item December has 31 days.
\end{enumerate}

\section*{Input}
The first line contains two integers $N$ and $d$ ($2 \leq N \leq 10^5$, $1 \leq d \leq 364$), the number of bears (including Björne),
and the number of days each bear sleeps.

The following $N$ lines contain the days when each bear goes into hibernation. They are given in the format \texttt{dd/mm}.

The first bear listed is Björne himself.

\section*{Output}
Print a line with an integer, the minimum number of bears that Björne must hire.

If it is not possible to guard the store every day, print \texttt{-1}.

\section*{Scoring}
Your solution will be tested on a set of test groups.
To earn points for a group, you must pass all test cases in that group.

\noindent
\begin{tabular}{| l | l | l |}
\hline
\textbf{Group} & \textbf{Points} & \textbf{Constraints} \\ \hline
$1$   & $30$       & $N \leq 15$ \\ \hline
$2$   & $40$       & $d \leq 182$ \\ \hline
$3$   & $30$       & No additional constraints. \\ \hline
\end{tabular}
